% 第 17 讲 消没定理的证明（译自 src/lecture17.txt）
\section*{第 17 讲\quad 消没定理的证明}

本讲的目标是证明：对非阿基米德局部域的扩张 $L/K$，设 Galois 群为 $G$，则
\[
\big[(L\otimes_K K^{\mathrm{unr}})^\times\big]_{tG}\simeq 0.
\]
回忆这蕴涵 $(L^\times)_{tG}\simeq \mathbb Z_{tG}[-2]$（即上同调局部类域论的主定理），后者又蕴涵 $G^{\mathrm{ab}}\simeq K^\times/NL^\times$。目前我们先假设 $L/K$ 完全分歧（一般情形到后面再归约），此时有 $L^{\mathrm{unr}}=L\otimes_K K^{\mathrm{unr}}$。上一次我们证明了：只需对每个素数 $\ell$ 与 $G$ 的每个 $\ell$-Sylow 子群 $G_\ell\subseteq G$ 证明
\[
\widehat H^0(G_\ell,L^{\mathrm{unr},\times})=0=\widehat H^1(G_\ell,L^{\mathrm{unr},\times}).
\]
注意，若令 $K':=L^{G_\ell}$，则 $L/K'$ 是 $G_\ell$-Galois 扩张。于是我们可以把 $K$ 换成 $K'$、把 $G$ 换成 $G_\ell$，从而不妨直接假设 $G$ 是 $\ell$-群（即 $\#G=\ell^n$，其中 $n$ 为某个整数）。

现在，上述两个等式中后一个不过是 Hilbert 定理 90（或习题集 7 第 3 题所给的推广）施用于扩张 $L^{\mathrm{unr}}/K^{\mathrm{unr}}$，因此还需证明前一个，即范映射
\[
N:L^{\mathrm{unr},\times}\to K^{\mathrm{unr},\times}
\]
是满射。

我们回忆 $\ell$-群的结构理论：

\textbf{命题 17.1.} 设 $G$ 是 $\ell$-群。则存在正规子群链
\[
1\lneq G_0\lneq\cdots\lneq G_m=G,
\]
使得 $G_{i+1}/G_i$ 对每个 $i$ 都是循环群。

\textit{证明.} 主要一步是证明：若 $G\neq 1$，则 $Z(G)\neq 1$（即 $G$ 的中心非平凡）。让 $G$ 通过伴随作用作用在自身之上，即对 $g,x\in G$ 令 $g\cdot x:=gxg^{-1}$。于是每个 $G$-轨道的大小或者为 $1$、或者被 $\ell$ 整除。由于
\[
\sum_{O\in G\text{-轨道}}\#O=\#G=\ell^n>1,
\]
而 $1$ 的 $G$-轨道大小为 $1$，故 $\ell$ 必须整除那些大小为 $1$ 的 $G$-轨道的个数，于是 $\ell\mid \#Z(G)$（特别地 $\#Z(G)\neq 0$）。% [原文作「hence #Z(G) ≠ 0」，按前文应为 ℓ | #Z(G)；照前文译]
然后取非平凡元素 $x\in Z(G)$，可见 $\langle x\rangle$ 是 $G$ 的正规子群，从而可对商群 $G/\langle x\rangle$ 作归纳即得结论。% [原文作「G/⟨x⟩ is a normal subgroup of G」，商群非子群；已按「⟨x⟩ 正规，对商群归纳」译]\hfill$\square$

于是由 Galois 理论，我们有一列相应的循环扩张
\[
L=L_m/L_{m-1}/\cdots/L_0=K.
\]
由于只需在这些子扩张的每一个上证明范映射是满射（满射的复合仍是满射），我们可以假设 $G$ 是循环群，比如说阶为 $n$。回忆 $N:L^\times\to K^\times$ 并非满射，因为我们已证明 $\#\widehat H^0(G,L^\times)=n$。现在，对每个 $m$，记 $K\subseteq K_m\subseteq K^{\mathrm{unr}}$ 为 $K$ 的 $m$ 次非分歧扩张。主要的一步如下：

\textbf{断言 17.2.} 设 $x\in K^\times$，并记 $L_m:=LK_m$。则对满足 $n\mid m$ 的 $m$，$x$ 属于 $N:L_m^\times\to K_m^\times$ 的像。% [原文只写「对每个 m」且未定义 L_m；按上下文（K_m 为 m 次非分歧扩张、下文 $N_{L_m/L}$）应取 L_m=L K_m。又因 $K^\times/NL^\times$ 的阶为 n，下面复合映射的「取 m 次幂为零」需 n | m（原文作取 n 次幂，见下），故补上此条件；下文只用到 m→∞ 的极限，无碍]

\textit{证明.} 注意 $K^\times/NL^\times=\mathcal O_K^\times/N\mathcal O_L^\times$。事实上，我们有通常的短正合列
\[
0\to \mathcal O_L^\times\to L^\times\xrightarrow{\ v\ }\mathbb Z\to 0,
\]
它给出正合列
\[
0=\widehat H^{-1}(\mathbb Z)\to \widehat H^0(\mathcal O_L^\times)\hookrightarrow \widehat H^0(L^\times)\to \widehat H^0(\mathbb Z)=\mathbb Z/n\mathbb Z,
\]
而最右端的映射为零，因为 $L/K$ 完全分歧（从而对一切 $y\in K^\times$ 有 $n\mid v(y)$）。于是得到同构 $\widehat H^0(L^\times)\simeq \widehat H^0(\mathcal O_L^\times)$，这正是我们的观察。

我们有交换图
\[
\begin{CD}
L^\times @>N>> K^\times @>>> K^\times/NL^\times @>>> 0\\
@VVV @VVV @VVV\\
L_m^\times @>N>> K_m^\times @>>> K_m^\times/NL_m^\times @>>> 0\\
@V{N_{L_m/L}}VV @V{N_{K_m/K}}VV @V{}VV\\
L^\times @>>> K^\times @>N>> K^\times/NL^\times @>>> 0.
\end{CD}
\]

现在，由诱导映射复合而成的映射 $K^\times/NL^\times\to K_m^\times/NL_m^\times\to K^\times/NL^\times$ 是取 $m$ 次幂（$m=[K_m:K]$），因此当 $n\mid m$ 时它是零映射。% [原文作「取 n 次幂」，但按上图复合（包含 $K^\times\hookrightarrow K_m^\times$ 接范映射 $N_{K_m/K}$）实为 $x\mapsto x^{[K_m:K]}=x^m$；又 $K^\times/NL^\times$ 的阶为 n，故须 n | m 才为零，原文此处的 n 与「对每个 m」皆有疏漏，已按数学含义订正] 我们想证明诱导映射
\[
N_{K_m/K}:K_m^\times/NL_m^\times\to K^\times/NL^\times
\]
是同构，而这蕴涵诱导映射 $K^\times/NL^\times\to K_m^\times/NL_m^\times$ 为零，从而证明断言。由断言 7.8(3)，即我们先前对 Herbrand 商的分析，两个群的阶都是 $n$，因此该映射单射当且仅当它满射。此外，由我们的观察，这等价于映射
\[
N_{K_m/K}:\mathcal O_{K_m}^\times/N\mathcal O_{L_m}^\times\to \mathcal O_K^\times/N\mathcal O_L^\times,
\]
而由于 $N:\mathcal O_{K_m}^\times\to \mathcal O_K^\times$ 是满射（其证明与断言 3.4 的证明完全相同），该映射也是满射，证毕。\hfill$\square$

再次设 $G$ 为 $n$ 阶循环群，我们只需证明 $N:L^{\mathrm{unr},\times}\to K^{\mathrm{unr},\times}$ 是满射。把断言施用于 $K_m^\times$，可见 $K_m^\times$ 的每个元素都是 $L_{m+m'}^\times$ 中某个元素的范数，因此
\[
N:\bigcup_m L_m^\times\to \bigcup_m K_m^\times
\]
是满射。接下来只需过渡到完备化。我们知道映射 $N:L^{\mathrm{unr},\times}\to K^{\mathrm{unr},\times}$ 的像包含 $\bigcup_m K_m^\times$，而后者稠密，故只需证明其像包含 $1$ 的一个开邻域。显然
\[
N(L^{\mathrm{unr},\times})\supseteq N(K^{\mathrm{unr},\times})=(K^{\mathrm{unr},\times})^n,
\]
而我们在习题集 1 第 1(a) 题中已经看到，$1+\mathfrak p_{K^{\mathrm{unr}}}^{2v(n)+1}$ 的每个元素都是 $K^{\mathrm{unr}}$ 中的 $n$ 次幂；这就是所要的开邻域。

最后，我们证明消没定理的一般情形，此时非阿基米德局部域的 $G$-扩张 $L/K$ 未必完全分歧。设 $L/L_0/K$ 为 $L$ 中 $K$ 的（唯一）极大非分歧扩张，于是 $L/L_0$ 完全分歧。令 $H:=\Gal(L/L_0)$，则 $L_0/K$ 是 Galois 扩张，群为 $G/H$。

\textbf{引理 17.3.} 设 $X$ 是 $G$-模的复形，并设有正合列
\[
1\to H\to G\to G/H\to 1.
\]
若
\[
X_{tH}\simeq 0\simeq (X^{hH})_{tG/H},
\]
则 $X_{tG}\simeq 0$。

注意，一般而言 $(X_{tH})_{tG/H}=X_{tG}$ 并不成立！例如，若 $H$ 是平凡群，则
\[
X_{tH}=\operatorname{hCoker}(N:X\to X)=0,
\]
其中此处 $N=\mathrm{id}_X$。

\textit{证明.} 由第一个条件，$X_{hH}\xrightarrow{qis}X^{hH}$；由习题集 6 第 3 题，有 $X^{hG}\simeq(X^{hH})^{hG/H}$ 与 $X_{hG}\simeq(X_{hH})_{hG/H}$。于是由第二个条件，
\[
(X^{hH})_{hG/H}\xrightarrow{qis}(X^{hH})^{hG/H}.
\]
容易验证这个拟同构由范映射给出（它是两个范映射的复合），这蕴涵
\[
X_{tG}=\operatorname{hCoker}(X^{hG}\to X^{hG})
\]
无环，如所欲证。\hfill$\square$

现在，我们想证明 $[(L\otimes_K K^{\mathrm{unr}})^\times]_{tG}\simeq 0$。回忆
\[
L\otimes_K K^{\mathrm{unr}}=L\otimes_{L_0}L_0\otimes_K K^{\mathrm{unr}}=L\otimes_{L_0}\prod_{L_0\hookrightarrow K^{\mathrm{unr}}}K^{\mathrm{unr}}=\prod_{L_0\hookrightarrow K^{\mathrm{unr}}}L\otimes_{L_0}K^{\mathrm{unr}}
\]
是典范地成立的（第二个同构由映射 $\alpha\otimes\beta\mapsto (i(\alpha)\beta)_i$ 给出，下标取遍嵌入 $i:L_0\hookrightarrow K^{\mathrm{unr}}$）。我们有
\[
\big[(L\otimes_K K^{\mathrm{unr}})^\times\big]_{tH}\simeq \prod_{L_0\hookrightarrow K^{\mathrm{unr}}}\big[(L\otimes_{L_0}K^{\mathrm{unr}})^\times\big]_{tH}\simeq \prod_{L_0\hookrightarrow K^{\mathrm{unr}}}\big[(L\otimes_{L_0}L_0^{\mathrm{unr}})^\times\big]_{tH}\simeq 0
\]
（用完全分歧情形，因为 $L_0/K$ 非分歧而 $L/L_0$ 完全分歧），这就建立了引理的第一个条件。为验证第二个条件，注意
\[
\prod_{L_0\hookrightarrow K^{\mathrm{unr}}}\big[(L\otimes_{L_0}K^{\mathrm{unr}})^\times\big]^{hH}=\prod_{L_0\hookrightarrow K^{\mathrm{unr}}}K^{\mathrm{unr},\times}\simeq K^{\mathrm{unr},\times}[G/H]
\]
作为 $G/H$-模成立（一旦固定一个嵌入 $L_0\hookrightarrow K^{\mathrm{unr}}$）。但正如习题集 7 第 1(e) 题所示，Tate 上同调对诱导模消没（因此上述等式无关紧要，我们只需要取遍这类嵌入的乘积来构造一个诱导 $G/H$-模）。于是引理 17.3 给出所要结论。
